\documentclass[11pt]{article}
\usepackage[a4paper,margin=24mm]{geometry}
\usepackage{amsmath,amssymb,amsthm,mathtools,mathrsfs}
\usepackage[hidelinks]{hyperref}
\setlength{\emergencystretch}{2em}
\allowdisplaybreaks
\newtheorem{theorem}{Theorem}
\newtheorem{lemma}[theorem]{Lemma}
\newcommand{\C}{\mathbb C}
\newcommand{\R}{\mathbb R}
\newcommand{\im}{\operatorname{im}}
\newcommand{\diag}{\operatorname{diag}}
\newcommand{\norm}[1]{\left\lVert#1\right\rVert}
\title{Native precision, marked detection and quantitative kernel leakage}
\author{Split-Zero research programme}
\date{19 September 2026}
\begin{document}\maketitle
This is the receiving proof for the supplied continuation
\cite{incoming}. It retains the original cutoff, full root polynomial,
period, unit phase, minimum metrics and marked scale. Its analytic
statements use the stated estimates of the preceding native proof;
their finite coefficient checks are not substituted for those estimates.
We prove the sharper precision consequences, the first observed marked
defect, the exact stack metric, and the certified heat depth. We also
derive an explicit positive lower bound for the original leakage Gram,
using the actual observation inverse and both action cross blocks.

\section{The precise analytic inputs and their receiving maps}
On the original simple-quartet domain let
\[
\begin{gathered}
 a\equiv1\pmod4,\quad q=(a+1)^2,\quad Q=(a-7)^2,\\
 \Delta=16a-48,\quad v=\operatorname{ord}_0E_A,
 \quad \epsilon\in\{0,1\},\quad s=1+\epsilon(a-1).
\end{gathered}
 \tag{NP1}
\]
Set $\ell=(s-1)/4$, $n=Q+\ell$, $m=\Delta-v$ and
$H=2\Delta-v-\ell$. These retain the original source order even
for the degree-$(a-8)$ lower root grid. Its four cutoffs are
$q-v-1,q-v,2q-v-1,2q-v$.

Use the actual equilibrium endpoints $u_e^2,v_e^2$ for
$\pi\sqrt{x}-2\log x$, satisfying $u_eK(\kappa)=2$,
$v_eE(\kappa)=4$, $\kappa^2=1-u_e^2/v_e^2$. Define
\[
\begin{gathered}
 w=(u_e+v_e)/2,\quad z_e=u_ev_e,\quad c_e=(v_e^2-u_e^2)/4,\\
 L=6\log w-2\log z_e-2,\quad
 \lambda_e=8-4\log w-2\log c_e,\quad F=-\lambda_e/2-3+L,\\
 C_B=9-8\log2+F,\quad C_\Gamma=2L-8\log2+4,\\
 C_\sigma=F+6-3L/2-6\log2+2\log\pi,\quad
 M_{-1}=(u_e^2+v_e^2-4u_ev_e)/(2u_e^2v_e^2),\\
 a_0=\log(4/\pi),\quad a_1=2(1-\log2)-\lambda_e/4,
 \quad b_1=\tfrac12\log c_e-\log2,\\
 C_\partial=2C_B-8a_1+4a_0,\quad
 K_2=C_B-4a_1+2b_1-2\log2,\quad
 \eta_h=(\gamma^2-\delta^2)/3.
\end{gathered}\tag{NP2}
\]
The quantitative input D17 of \cite{native} is
\[
\begin{aligned}
\mathcal C^-_{a,s}={}&C_Bn^2+4n(Ha_1-ma_0)
+\tfrac12m^2\log(n/m)+2H^2b_1\\
&-(2a_0-\tfrac34)m^2-2\ell^2\log2+C_\sigma n
-2\eta_h Q(Q-1)M_{-1}/n+O_{h,v}(a^{3/2}+a\log(a+2)).
\end{aligned}\tag{NP3}
\]
Its proof sums the low rows with their full endpoint weights, including
$\sum c_rr=m^2$ and
$\sum c_rr\log r=m^2\log m-m^2/2+O(\log(m+2))$.
The expansion $\psi(t)=a_0+(t/4)\log t+(a_0-1/4)t+O(t^{3/2})$
gives the $-3/4$ in (NP3): $-1/2$ from the linear summand and
$-1/4$ from this weighted logarithmic sum. No interval is omitted.

\begin{theorem}
Let
\[
 Z(v)=1024(a_1-2b_1)-176C_\partial+4v(a_1-a_0)-192+512\log2,
 \qquad Z_\Delta=4C_\partial+16(2b_1-a_1).
 \tag{NP4}
\]
The expanded lower covariance satisfies
\[
\begin{aligned}
\mathcal C^-_{a,s}={}&C_Bq^2+[-16C_\partial+\epsilon C_\Gamma/4]aq
+128q\log a\\
&+q[C_\sigma-2\eta_hM_{-1}-\epsilon C_\Gamma/4
+\epsilon K_2/16-Z(v)-\epsilon Z_\Delta]\\
&+O_{h,v}(a^{3/2}+a\log(a+2)).
\end{aligned}\tag{NP5}
\]
The complete original correlated native sum obeys
\[
 \mathcal N_{a,s}=16C_\partial aq-128q\log a
    +[Z(v)+\epsilon Z_\Delta]q
       +O_{h,v}(a^{3/2}+a\log(a+2)). \tag{NP6}
\]
After the displayed terms, each remainder is $o(q/\log q)$; in
particular its next $q/\log q$ coefficient is zero. Also
$\mathcal N_{a,a}-\mathcal N_{a,1}=Z_\Delta q+o(q/\log q)$.
\end{theorem}
\begin{proof}
Substitute $n=a^2+(-14+\epsilon/4)a+49-\epsilon/4$ into (NP3).
All polynomial terms of degrees four, three and two yield (NP5)
by expansion. In the logarithmic term,
$\log(n/m)=\log a-\log16+O_{h,v}(a^{-1})$,
so its discarded part is $O(a)$; changing its degree-two coefficient
from $a^2$ to $q$ costs $O(a\log a)$. Further,
$Q(Q-1)/n=q+O_{h,v}(a)$ and $n=q+O(a)$. These estimates account
for every nonpolynomial replacement. The supplied exact coefficient
diagnostics independently check the three polynomial coefficients
for both $\epsilon=0,1$.

The original receiving identity is
$\mathcal B_a^{(s)}=\mathcal C^-_{a,s}+\mathcal N_{a,s}$, where
$\mathcal N=\mathcal Re^C+\mathcal T+F_L+\Xi$ retains all four
specified terms. The upper Gamma baseline is
\[
 \mathcal B_a^{(s)}=C_Bq^2+\epsilon C_\Gamma aq/4
 +q[C_\sigma-2\eta_hM_{-1}-\epsilon C_\Gamma/4+\epsilon K_2/16]
 +O_h(a\log(a+2)).
\]
Its order-one remainder is $O_h(\log q)$ in the included absolute
baseline proof; D20's order conversion costs $O_h(a\log q)$.
Subtracting proves (NP6). Finally
$(a^{3/2}+a\log(a+2))\log q/q\to0$. The claims concern this
complete correlated sum and covariance, not their individual
restriction or quotient determinants.
\end{proof}

At fixed original multiplicity $m_0$, retain
$q=[1+k(m_0-1)](k+1)^2$, $\ell=(k-1)/4$ and the original
arithmetic density. D20 gives the exact source-order comparison
through its stated error. With the earlier arithmetic result SR/AG
as used in FW29--30 \cite{FW}, subtraction yields
\[
\begin{aligned}
\delta_{k,k}^{\rm ar}={}&-(\ell+4m_0-2)C_\Gamma q
+\frac{\log2}{\pi}I_hk^2-\ell^2K_2
+\frac{C_\Gamma q}{2(\log q+c_\zeta)}+o_h(q/\log q),\\
 I_h={}&\int_0^{\pi/2}[(\log M_h)'(t)]^2dt,
\quad M_h(t)=\int_{\R}e^{ty}w_h(y)dy,
\quad c_\zeta=2\gamma_{\rm E}-\log(2\pi).
\end{aligned}\tag{NP7}
\]
Here $w_h(y)=|2\xi(1/2+iy)/h(1/2+iy)|^2/(2\pi)$ and $h$ is
the complete fixed quartet polynomial. The factor $k^2$ remains
literal. For $m_0>1$ an $o(q/\log q)$ remainder does not determine
every smaller $k^2$-scale term separately. The proof is subtraction of
$\ell qC_\Gamma+\ell^2K_2+O_h(k\log q)$ from the order-one
correction; $k(\log q)^2/q\to0$ proves the error claim. The exact
complete-action receiver is
\[
 \mathcal B_k[w_k]+\delta_{k,k}^{\rm ar}
       =\mathcal B_k[\sigma]+\delta_{k,1}^{\rm ar}. \tag{NP8}
\]
The larger conversion terms cancel in this identity.

For the sign in D21, set $r=u_e/v_e$ and $\kappa=\sqrt{1-r^2}$.
Differentiating the elliptic integrals using Carlson's formulas
\cite[19.4.1--2]{Carlson} gives exactly
\[
 \frac{d}{dr}\frac{rK(\kappa)}{E(\kappa)}
 =\frac{(K-E)(E-r^2K)}{(1-r^2)E^2}>0. \tag{NP9}
\]
The first factor is positive by the integral of
$\kappa^2\sin^2t/\sqrt{1-\kappa^2\sin^2t}$; the second by the
integral of $\kappa^2\cos^2t/\sqrt{1-\kappa^2\sin^2t}$.
For completeness, the substitution $x=\tan t$ bounds
$K=\int_0^\infty[(1+x^2)(1+r^2x^2)]^{-1/2}dx$
by $1+\log(1/r)+1$, using the intervals $[0,1]$, $[1,1/r]$,
and $[1/r,\infty)$. Also $E\ge\int_0^{\pi/2}\cos t\,dt=1$.
At $r=1/10$, $rK/E<(2+3)/10=1/2$, so the actual solution
of $rK/E=1/2$ exceeds $1/10$. Since
$1<(1+r)^2/(4r)<121/40$, the exact identity
$K_2=2\log(w^2/z_e)-4\log2$ yields
$-4\log2<K_2<2\log(121/160)<0$.

\section{The first visible jet and the original marked recurrence}
For this section use the original simple quartet and $k\ge9$,
$k\equiv1\pmod4$. Retain its actual invertible period matrix
$H(\varpi)$, and put
\[
 x(z)=H(\varpi)(e^{2(\delta+i\gamma)z},e^{2\delta z},
 e^{2i\gamma z},1)^T,\quad b_*=1/2-\delta-i\gamma.
 \tag{NP10}
\]
The original observation pairs $P(S)$ with
$P(\partial_z)[e^{kb_*z}f(x(z))]_{z=0}$, for invariant homogeneous
$f$ of degree $k$ with coordinate weights $1,2,3,4$ modulo five.
Define $j_* =\min\{j:\Lambda S^j\ne0\}$.

\begin{theorem}
If $x(0)$ has two active coordinates then $j_*=0$. If only $x_i(0)$
is active, write $z_j=x_j/x_i$, $\nu_j=\operatorname{ord}_0z_j$
and $c_j=j-i\pmod5$. Then
\[
 j_* =\min_{\substack{\alpha\ge0,\ |\alpha|\le4\\
             \sum_{j\ne i}c_j\alpha_j=-ki\ (\mathrm{mod}\ 5)}}
                  \sum_{j\ne i}\nu_j\alpha_j\le4. \tag{NP11}
\]
In this case $j_*=0$ exactly when $5\mid k$.
\end{theorem}
\begin{proof}
If coordinates $i,j$ are active, choose $0\le b\le4$ with
$ki+b(j-i)=0\pmod5$. The monomial $x_i^{k-b}x_j^b$ is invariant
and nonzero at zero. In the one-active case, each inactive coordinate
is a nonzero four-exponential sum: its row in $H$ is nonzero. The
first four exponential derivatives give an invertible Vandermonde,
so $1\le\nu_j\le3$. At least one equals one, since otherwise
$x'(0)$ is proportional to $x(0)$, which invertibility of $H$ would
make the distinct exponent vector proportional to $(1,1,1,1)$.

Divide each invariant monomial by the common unit
$e^{kb_*z}x_i(z)^k$. Its charge is exactly that in (NP11), and its
order is the indicated sum. A linear combination cannot have smaller
order than every summand. Conversely each candidate of degree at most
four homogenizes to degree $k$ and has that exact order. A power at
most four of an order-one ratio realizes every charge; any monomial
of total degree at least five has order at least five and cannot
improve the minimum. The charge-zero constant realizes order zero;
all other charges have positive order. This proves every assertion.
\end{proof}

Let $E=\C[S]/(\chi)$, $\deg\chi=q$, in the original increasing
power basis. Denote multiplication by $S$ by $A$, and let
$B=e_0e_{q-1}^T$. Retain the original marked scale $u\ne0$ and
the connection
\[
 A(t)=A+tB,\quad U'(t)=A(t)U(t)/u,\quad U(0)=I,
 \quad\Lambda_h(t)=\Lambda U(t)^{-1}. \tag{NP12}
\]
For the actual terminal class $v_\lambda=b_\lambda\chi/(S-\lambda)$,
one has $Av_\lambda=\lambda v_\lambda$ and $Bv_\lambda=b_\lambda1$.
Thus, with $Y(t)=e^{\lambda t/u}\Lambda_h(t)v_\lambda-\Lambda v_\lambda$,
\[
 Y'(t)=-\frac{b_\lambda}{u}t\,e^{\lambda t/u}\Lambda_h(t)1.
 \tag{NP13}
\]
This follows by differentiating $U^{-1}$ on its correct side:
$(U^{-1})'=-U^{-1}A(t)/u$.

To calculate the first coefficient without an unjustified commutation,
write $U^{-1}=\sum W_nt^n$. Then
$(n+1)W_{n+1}=-(W_nA+W_{n-1}B)/u$, $W_{-1}=0$.
In any word contributing to $W_n$, the letter $A$ has degree one
and $B$ degree two. A word containing $B$ kills $1$ unless its
rightmost $B$ has at least $q-1$ letters $A$ to its right, since
$BA^j1=0$ for $j<q-1$. Hence
$W_n1=(-A/u)^n1/n!$ for $n<q+1$. Since $j_*\le4<q$,
(NP13) proves
\[
 Y^{(j_*+2)}(0)=b_\lambda(j_*+1)(-u)^{-(j_*+1)}
                    \Lambda S^{j_*}\ne0,\qquad
 Y^{(j)}(0)=0\quad(j<j_*+2). \tag{NP14}
\]
The first marked/arithmetic discrepancy is therefore at order two on
the two-active stratum, and at an exactly determined order at most six
at every admitted period. This computes the stated holomorphic
continuation, with full $b_\lambda$ and $u$; it does not substitute a
new sheaf morphism for that continuation.

\section{Observation stacks in their actual quotient metric}
Fix an original cutoff $N$, its positive form $G=G_N$, and the
surjective observation $\Lambda:E\to B_0=\im\Lambda$. Define
\[
 Q=(\Lambda G^{-1}\Lambda^*)^{-1},\quad
 L=G^{-1}\Lambda^*Q,\quad K=\ker\Lambda,\quad
 I:\C^r\longrightarrow K,\quad H_K=I^*GI. \tag{NP15}
\]
$I$ is the original fixed full-column kernel frame. Thus $\Lambda L=I$
on $B_0$, $I^*GL=0$ and $L^*GL=Q$. Set
$J_K=H_K^{-1}I^*G$, so $IJ_K+L\Lambda=I_E$.
The literal action blocks are
\[
 T=J_KAI,\quad C=J_KAL,\quad B=\Lambda AI,
 \quad D=\Lambda AL,\quad
 A[I,L]=[I,L]\begin{pmatrix}T&C\\B&D\end{pmatrix}.
 \tag{NP16}
\]
All norms below use $H_K$ on kernel coordinates and $Q$ on outputs.

The incoming O1--9 theorem \cite{native} proves $K_\infty=0$ on its
explicit original period locus, with $r=8k-16$, and injectivity of
$\mathcal O_r=(\Lambda,\Lambda A,\ldots,\Lambda A^r)^T$.
Its proof is the decreasing-kernel argument: a plateau is invariant
under $A$, hence zero; at most $r$ strict drops start from $K$.
Consequently
\[
 \mathcal L_r=\sum_{j=0}^{r-1}(T^j)^*B^*QB T^j\succ0.
 \tag{NP17}
\]
Indeed if all $BT^jx=0$ for $j<r$, the characteristic polynomial
of $T$ gives the same for every $j$. The block recurrence then gives
$A^jIx=IT^jx$, so $Ix\in K_\infty$ and $x=0$.

The full stack metric is
\[
 \mathsf H_{d,N}=\sum_{j=0}^d(A^j)^*\Lambda^*Q\Lambda A^j,
 \quad v_\lambda^*\mathsf H_{d,N}v_\lambda
  =\left(\sum_{j=0}^d|\lambda|^{2j}\right)E_{N,\lambda}(1-\theta_N),
 \tag{NP18}
\]
where $E_{N,\lambda}=v_\lambda^*G_Nv_\lambda$ and
$\theta_N=\|J_Kv_\lambda\|_{H_K}^2/E_{N,\lambda}$.
The equality is substitution of the actual eigenvalue equation into
each term. For fixed $d$ the geometric factor cancels in the original
four-cutoff return; FW38--40 \cite{FW} therefore yields
\[
 C_\partial q+
 \log\frac{(1-\theta_{q-1})(1-\theta_q)}
               {(1-\theta_{2q-1})(1-\theta_{2q})}
       +O_h(k+\log q). \tag{NP19}
\]
The stack retains the original angle term by an exact identity.

\section{A quantitative lower bound for every kernel direction}
We now extend positivity (NP17) by a bound expressed in actual finite
matrices. This is valid on the just-stated simple-packet observability
locus and at every original cutoff. It does not replace that cutoff.

Let the distinct arithmetic roots be $\lambda_1,\ldots,\lambda_q$,
with original Lagrange idempotents
$e_\lambda(S)=\prod_{\mu\ne\lambda}(S-\mu)/(\lambda-\mu)
=\sum_{j=0}^{q-1}e_{\lambda,j}S^j$. Write
$b_\lambda^{\rm obs}=\Lambda e_\lambda\ne0$ and define
\[
 \vartheta_\lambda(y)=
 \frac{(b_\lambda^{\rm obs})^*Qy}{
             (b_\lambda^{\rm obs})^*Qb_\lambda^{\rm obs}},
 \quad\mathcal R_N(y_0,\ldots,y_{q-1})
 =\sum_\lambda e_\lambda\sum_{j=0}^{q-1}e_{\lambda,j}
                                      \vartheta_\lambda(y_j).
 \tag{NP20}
\]
For the original stack, its $\lambda$ coordinate is
$\vartheta_\lambda(\Lambda e_\lambda(A)x)=x_\lambda$, so
$\mathcal R_N\mathcal O_{q-1}=I_E$. This is O10 with its
functional explicitly placed in the original quotient metric $Q$;
no metric equivalence is silently assumed. By Cauchy--Schwarz,
\[
 \|\mathcal R_N\|_{Q^{\oplus q}\to G}\le C_{R,N}:=
 \sum_\lambda\frac{\|e_\lambda\|_G
                 (\sum_j|e_{\lambda,j}|^2)^{1/2}}
                        {\|b_\lambda^{\rm obs}\|_Q}<\infty.
 \tag{NP21}
\]
This coefficient contains the original root differences and metric.

For $x\in\C^r$ write
$y_j=\Lambda A^jIx$, $x_j=J_KA^jIx$. The exact block recurrences
are $x_{j+1}=Tx_j+Cy_j$, $y_{j+1}=Bx_j+Dy_j$,
with $x_0=x,y_0=0$. Elimination of $x_j$ gives
\[
 y_{j+1}=BT^jx+Dy_j+
                 \sum_{l=0}^{j-1}BT^{j-1-l}Cy_l.
 \tag{NP22}
\]
Define a $(q-1)$ by $(q-1)$ block lower-triangular matrix $W_N$ by
the following recursion. For input $z_0,\ldots,z_{q-2}$ set $y_0=0$
and $y_{j+1}=z_j+Dy_j+\sum_{l=0}^{j-1}BT^{j-1-l}Cy_l$.
Then $W_Nz=(y_1,\ldots,y_{q-1})$. Its diagonal blocks are identities;
the displayed reverse formula gives its exact inverse. In particular,
both action cross blocks $B,C$ and $D$ are retained.

Let $\chi_T(t)=\det(tI-T)$, and compute the actual polynomial
remainders
$t^j\bmod\chi_T=\sum_{h=0}^{r-1}c_{jh}t^h$, $0\le j\le q-2$.
Put $C_T=(c_{jh})$. The first $r$ rows form the identity. Applying
the characteristic identity to $T$ proves
\[
 (BT^jx)_{j=0}^{q-2}=(C_T\otimes I_{B_0})(BT^hx)_{h=0}^{r-1},
 \quad
 \mathcal O_{q-1}Ix=
 (0,\ W_N(C_T\otimes I_{B_0})(BT^hx)_{h=0}^{r-1}).
 \tag{NP23}
\]
Define $C_{W,N}=\|W_N\|_{Q^{\oplus(q-1)}}$ and
$C_{T,N}=\|C_T\|_2$. These are norms of the displayed finite
matrices, not constants of a different stack. Then
\[
 \boxed{\quad
 \mathcal L_r\succeq
    (C_{R,N}C_{W,N}C_{T,N})^{-2}H_K.
 \quad} \tag{NP24}
\]
To prove it, apply (NP20) to (NP23), take the original $G$ norm,
and then use (NP21) and the two matrix norms. The square of
$\|(BT^hx)_{h<r}\|_{Q^{\oplus r}}$ is exactly
$x^*\mathcal L_rx$; $\|Ix\|_G^2=x^*H_Kx$. Rearrangement gives
(NP24). It supplies a strictly positive finite lower bound at the
original cutoff. Its dependence on the actual $Q$ and action blocks
is displayed. No small asymptotic action allowance is inferred from
this possibly large recovery cost.

\section{The exact heat guard and its evaluated degree cost}
Retain HR1--11 \cite{heat}, including the original complete
prequotient constant $c_N\ge1$ and the same vectors and observation.
Let $0<p_{1,N},p_{2,N}\le1$ be the actual squared source-correlation
fractions from FW5, and $p_{*,N}=\min(p_{1,N},p_{2,N})$.
For $0<\zeta<1$, the certificate
\[
 \eta_{J,N}=(1+2^{-(J+1)/2}c_N)^{2k}-1
       \le\zeta\sqrt{p_{*,N}}
 \tag{NP25}
\]
is equivalent, by monotonicity, to
\[
 J+1\ge T_N:=2\log_2c_N-2\log_2
 \left[(1+\zeta\sqrt{p_{*,N}})^{1/(2k)}-1\right].
 \tag{NP26}
\]
The least nonnegative integer depth is $J_N^*=\lceil T_N\rceil-1$.
Indeed $c_N\ge1$ and the bracket is less than one; for the original
$k\ge9$, $T_N>1$. There is therefore no hidden cutoff at depth zero.

For $0<x\le\zeta$ put $b=\log(1+x)/(2k)$. The elementary bounds
$x/(1+\zeta)\le\log(1+x)\le x$ and
$b\le e^b-1\le(1+\zeta)^{1/2}b$ show, uniformly even as
$p_{*,N}\downarrow0$, that
\[
 J_N^*+1-2\log_2c_N
 =2\log_2(2k/\zeta)-\log_2p_{*,N}+O_\zeta(1).
 \tag{NP27}
\]
Here the equality concerns the least certified depth, not an arbitrary
larger $J$. The rounding error is less than one.

Let the original endpoint weights be $w_0=w_q=1$, $w_r=2$
otherwise. They sum to $2q$. FW27--29 supplies
$\mathscr A_1=-\sum w_r\log p_{1,q-1+r}=C_Bq^2+O_h(q\log q)$
and $\mathscr A_2=O_h(q\log q)$. Since
\[
 -\log p_{1,N}\le-\log p_{*,N}
       \le-\log p_{1,N}-\log p_{2,N},
\]
weighted summation of (NP27) proves
\[
 \frac1{2q}\sum_{r=0}^qw_r
 [J_{q-1+r}^*+1-2\log_2c_{q-1+r}]
 =\frac{C_B}{2\log2}q+O_{h,\zeta}(\log q).
 \tag{NP28}
\]
On the simple-quartet domain FW48 at $r=0$ gives
$-\log p_{1,q-1}=2a_0q+O_h(k+\log q)$ and
$-\log p_{2,q-1}=O_h(\log q)$. The same two-sided inequality yields
\[
 J_{q-1}^*+1-2\log_2c_{q-1}
 =\frac{2\log(4/\pi)}{\log2}q+O_{h,\zeta}(k+\log q).
 \tag{NP29}
\]
(NP28) uses the all-fixed-multiplicity source products; (NP29) uses
the stated simple-packet pointwise estimate. These are costs of this
particular heat certificate. They do not bound every possible numerical
method or evaluate the separate growth of $c_N$.

\section{Return to the now-proved coprime proper-source domain}
The latest published NC1--27 proof \cite{NC} establishes the
coprime stratum required in Section 5 of the incoming continuation:
outside a specified finite exceptional set of admitted periods,
$E_k^{\rm prop}=0$, with a sufficient original finite-period radius.
On its intersection with the two-active jet locus,
$r_I=8k-32-v$. At the same proper-source cutoffs the source form is
exactly $Q_N^p=Z_R^*J_0^*D_{k,N-v}^{(s_0)}J_0Z_R$.

In the full target product let $F_R$ be the actual multiplier lift,
$J_*$ the entire target-invariant direction matrix and $\mathcal D_N^p$
the unchanged target form. Define
\[
 A_N=F_R^*\mathcal D_N^pF_R,\quad
 C_N=J_*^*\mathcal D_N^pJ_*,\quad
 B_N=J_*^*\mathcal D_N^pF_R,\quad T_N=C_N^{-1/2}B_NA_N^{-1/2}.
 \tag{NP30}
\]
The symbol $T_N$ here denotes this displayed target cross map, not the
real heat threshold in (NP26). Completing the square in the entire
target fiber gives
$H_N^p=A_N-B_N^*C_N^{-1}B_N
=A_N^{1/2}(I-T_N^*T_N)A_N^{1/2}$.
Thus the original proper-source return is exactly
\[
 \mathcal M^p=\mathcal R_{\rm end}\log\det[(Q_N^p)^{-1/2}A_N(Q_N^p)^{-1/2}]
       +\mathcal R_{\rm end}\log\det(I-T_N^*T_N). \tag{NP31}
\]
This determinant identity follows by multiplying determinants, not by
asserting the matrices commute. The actual positive target minimum
gives $I-T_N^*T_N\succ0$. None of the covariance, observation-stack
or heat-depth results evaluates its eigenvalues. They now supply
sharper original inputs and a finite positive leakage bound for the
continuing calculation of this same target metric.

\begin{thebibliography}{9}
\bibitem{incoming} Supplied web continuation,
\emph{The new continuation removes more uncertainty than its concluding
summary states}, 19 September 2026, equations (1)--(23).
The complete original Markdown is retained with this edition.
\bibitem{native} \emph{Programme Native Continuation}, four complete
notes supplied 19 September 2026, D1--25, revised J1--6 including
J3a--b, E1--4, O1--12. Published with the NC edition below.
\bibitem{FW} \emph{Full-window source products and the retained class},
FW1--52, 19 September 2026, public commit
\texttt{3c44de632b2c5ead13b23619d646803d2bbffc19};
full proof \texttt{FULL\_WINDOW\_SOURCE\_PRODUCTS.tex} included.
\bibitem{heat} \emph{Heat approximation errors for the original complex
observation responses}, HR1--11, and the full original heat metric
proof HM1--30, published in commit
\texttt{607728ee3c35be98720d548794727670b1946c91}.
\bibitem{NC} \emph{Vanishing of the original proper-source kernel through
conductor coprimality}, NC1--27, 19 September 2026, commit
\texttt{12dd3645229bfebbf1cb8f5968ecffd4dfff1e9f}.
\bibitem{Carlson} Bille C. Carlson, \emph{Elliptic Integrals}, Chapter 19
of the NIST Digital Library of Mathematical Functions,
\href{https://dlmf.nist.gov/19.4}{Section 19.4}, equations 19.4.1--2.
The derivative in (NP9) is derived from these formulas; its sign
and the subsequent inequalities are proved from the defining integrals.
\end{thebibliography}
\end{document}
